\chapter{Introduction}


\section{Quantum control for near-term quantum computing}
\label{subsec:intro_quantum_control_nisq}

Quantum computing is ultimately a control problem as much as it is a computational one. At the abstract level, a quantum algorithm is formulated as a sequence of logical unitary operations acting on qubits; at the hardware level, however, those operations must be realized through continuous-time dynamics generated by experimentally accessible Hamiltonians. The practical power of a quantum device therefore depends not only on the existence of an algorithmic advantage, but also on the ability to engineer control protocols that implement the required gates with sufficiently high accuracy, sufficiently short runtime, and sufficient robustness against physical imperfections \cite{Niu2019,Koch2022}. This issue is especially acute in near-term and pre-fault-tolerant quantum processors, where limited coherence times, calibration uncertainty, waveform distortions, and multilevel effects strongly constrain the depth and reliability of executable circuits \cite{Krantz2019,Kjaergaard2020}. In this regime, quantum control is not a secondary layer of fine-tuning but a central bottleneck linking the abstract circuit model to the actual computational capacity of the hardware. For this reason, understanding how gate fidelity, total runtime, leakage, and robustness compete with one another is essential for assessing the practical viability of analog gate design, as will be discussed in Sections~\ref{subsec:fidelity_runtime_metrics}, \ref{subsec:computational_subspace_leakage}, \ref{subsec:speed_leakage_tradeoff}, and \ref{subsec:why_robustness_is_difficult}. Within this broader context, the reference work analyzed in this thesis is particularly relevant because it proposes a unified framework in which these competing objectives are encoded into a single optimization problem and solved through reinforcement learning on a realistic superconducting-qubit platform \cite{Niu2019}; its main ideas and results will be introduced in Sections~\ref{subsec:ufo_cost_function} and \ref{subsec:main_findings_original_paper}.

\section{Direct analog gate implementation versus compiled gate synthesis}
\label{subsec:intro_analog_vs_synthesis}

A central question in quantum gate design is whether a target unitary should be implemented through a compiled sequence of elementary gates or realized directly through a single analog control protocol. In the standard circuit-model view, universality is achieved by expressing a desired operation as a product of gates drawn from a finite universal set, an approach that is indispensable in fault-tolerant architectures and remains the dominant abstraction in quantum algorithm design. However, on near-term hardware this decomposition is not neutral from a physical perspective: every additional elementary gate consumes runtime, increases exposure to decoherence and control noise, and may introduce further leakage or calibration error. As a result, a circuit that is optimal or near-optimal at the logical level may still be inefficient at the hardware level once translated into pulses and executed on a concrete device \cite{Niu2019,NielsenChuang2010,Koch2022}. For pre-fault-tolerant platforms, this observation motivates a different perspective in which one seeks to implement the target unitary directly through continuous-time Hamiltonian evolution, thereby exploiting the native interaction structure of the hardware rather than forcing the operation through a fixed discrete gate alphabet.

This distinction between compiled synthesis and direct analog realization is particularly relevant for the family of two-qubit gates considered in the reference paper. There, the authors explicitly argue that unrestricted analog control can reduce the resource overhead associated with long synthesized circuits, and they benchmark their learned controls against optimal gate-synthesis counterparts precisely to quantify this hardware-efficiency advantage \cite{Niu2019}. The significance of this comparison is not only algorithmic but physical: if a target unitary is well aligned with the natural couplings of the device, then a pulse-level implementation may achieve the same logical transformation with substantially shorter runtime than a decomposition into standard one- and two-qubit primitives. This idea will be developed more systematically in Sections~\ref{subsec:universal_gates_analog_control}, \ref{subsec:gmon_architecture_motivation}, and \ref{subsec:target_gate_family}, and will reappear in the numerical part of the thesis when the reproduced analog controls are compared against gate-synthesis baselines in terms of total gate time; see in particular Sections~\ref{subsec:baselines_comparison_strategy} and \ref{subsec:main_findings_original_paper}. In this sense, the present thesis is motivated not only by the general problem of gate optimization, but also by the more specific question of when direct analog quantum control can outperform compiled gate implementations on realistic superconducting hardware.

\section{Practical limitations: leakage, stochastic errors, and robustness}
\label{subsec:intro_practical_limitations}

The promise of direct analog control is accompanied by equally important practical difficulties. Real quantum devices are not ideal few-level systems manipulated by perfectly calibrated controls, but open and imperfect physical platforms subject to multilevel dynamics, stochastic fluctuations, and partial model uncertainty. In superconducting circuits, in particular, the computational qubit subspace is embedded in a larger Hilbert space associated with higher oscillator levels, so fast or strongly modulated controls can drive population outside the logical sector and generate leakage errors. At the same time, the implemented control waveforms are affected by amplitude noise, parameter drift, transfer-function distortions, and calibration mismatch, all of which can substantially degrade the quality of the final gate if they are not accounted for during optimization \cite{Niu2019,Krantz2019,Kjaergaard2020,Koch2022}. These non-idealities imply that achieving high nominal fidelity alone is not sufficient: a practically useful control protocol must also remain stable under realistic perturbations and must suppress leakage while preserving short runtime.

The reference paper is motivated precisely by the lack of a unified framework capable of treating these competing requirements simultaneously. Its starting point is that leakage and stochastic control errors should not be regarded as secondary corrections to an otherwise solved nominal-control problem, but as structural features of the optimization landscape itself \cite{Niu2019}. This observation is central for the present thesis as well. Indeed, much of the background developed later, from the multilevel description of superconducting qubits to the definition of computational and leakage subspaces, from the speed--leakage tradeoff to the discussion of robustness under uncertain perturbations, serves to make this point precise; see Sections~\ref{subsec:superconducting_qubits_multilevel}, \ref{subsec:computational_subspace_leakage}, \ref{subsec:speed_leakage_tradeoff}, and \ref{subsec:why_robustness_is_difficult}. It is exactly because these physical limitations are intertwined that the paper introduces an analytic leakage treatment, a multi-objective cost function, and a stochastic training environment within a reinforcement-learning framework. These ingredients, which will be discussed in Sections~\ref{subsec:leakage_bound_penalty}, \ref{subsec:ufo_cost_function}, and \ref{subsec:stochastic_training_environment}, provide the conceptual basis for the experiments of this thesis, including the reproduction of the original robustness results and the out-of-distribution robustness tests proposed in the extension part of the work.

\section{The reference framework: \textit{Universal Quantum Control through Deep Reinforcement Learning}}
\label{subsec:intro_reference_framework}

The present thesis takes as its main point of departure the framework introduced in \textit{Universal Quantum Control through Deep Reinforcement Learning} by Niu \textit{et al.} \cite{Niu2019}. The central idea of that work is that realistic quantum gate design should be formulated as a genuinely multi-objective control problem, in which gate accuracy, runtime, leakage suppression, and robustness against stochastic control errors are optimized jointly rather than sequentially or in isolation. Instead of treating leakage and control noise as secondary effects to be evaluated only after a nominal pulse has been found, the paper incorporates them directly into the optimization target through a unified control objective. This is a conceptually important step, because it shifts the emphasis from finding merely high-fidelity nominal solutions to finding controls that remain physically meaningful under the non-ideal conditions relevant to superconducting quantum hardware.

To realize this program, the paper combines three main ingredients. First, it introduces an analytic treatment of leakage based on the time-dependent Schrieffer--Wolff transformation, allowing the dominant leakage mechanisms to be organized into a perturbative bound that can be used during optimization. Second, it defines the so-called Universal cost Function control Optimization (UFO) objective, which augments gate infidelity with explicit penalty terms for leakage, boundary-condition violation, and total runtime. Third, it casts the resulting control problem into a continuous-action reinforcement-learning framework and trains the control policy with trust-region policy optimization in a stochastic environment designed to mimic experimentally relevant control fluctuations \cite{Niu2019}. The framework is implemented on a two-qubit superconducting gmon architecture and benchmarked on a structured family of two-qubit gates that are relevant both from the viewpoint of hardware-native interactions and from applications in quantum simulation.

This reference framework is particularly suitable as the foundation of the present thesis for two reasons. On the one hand, it is physically rich: it does not rely on an idealized two-level approximation, but explicitly addresses multilevel dynamics, leakage channels, and stochastic perturbations in a realistic superconducting-qubit setting. On the other hand, it is methodologically modular: its physical model, leakage analysis, cost function, reinforcement-learning formulation, and evaluation strategy can each be reproduced, tested, and extended independently. For this reason, the work developed in this thesis is organized around a two-stage objective. The first is to reconstruct and validate the main ingredients and key numerical results of the original framework. The second is to extend that framework along directions not fully explored in the original paper, including out-of-distribution robustness tests, comparisons against more standard quantum optimal-control baselines, and a systematic study of the role played by the different terms of the UFO objective. The technical content of the framework itself will be developed in detail in Sections~\ref{subsec:gmon_effective_hamiltonian}--\ref{subsec:main_findings_original_paper}, while its reproduction and extension constitute the core of the empirical chapters of this thesis.


\section{Objectives and scope of this thesis}
\label{subsec:intro_objectives_scope}

The aim of this thesis is to study, reproduce, and extend a recent reinforcement-learning-based framework for robust analog quantum control in superconducting qubits. More specifically, the work takes the control strategy proposed in Ref.~\cite{Niu2019} as its main methodological reference and investigates both its physical content and its empirical performance on a realistic multilevel control problem. The scope of the thesis is therefore twofold. On the one hand, it is a \emph{reproduction study}: the physical model, the leakage-aware control objective, the reinforcement-learning formulation, and the principal numerical benchmarks of the original paper are reconstructed in a consistent simulation pipeline. On the other hand, it is an \emph{extension study}: once the original framework has been validated, additional experiments are carried out to test its robustness beyond the specific setting considered in the reference work.

The first objective is thus to reproduce the main results of the original paper as faithfully as possible. In practical terms, this means implementing the gmon control model, the UFO-inspired optimization objective, and the learning environment described in Ref.~\cite{Niu2019}, and then verifying whether the principal reported trends can be recovered. The core reproduction targets are the comparison of analog gate times against synthesized-gate baselines, the behavior of average gate fidelity as a function of control-noise strength, and the corresponding fidelity variance under stochastic perturbations. These results are particularly important because they are the main empirical support for the central claims of the paper, namely that noisy-environment reinforcement learning can improve both robustness and hardware efficiency relative to simpler baselines \cite{Niu2019}. In the structure of this thesis, this first objective motivates the material developed in Sections~\ref{subsec:gmon_architecture_motivation}--\ref{subsec:main_findings_original_paper} and culminates in the reproduction experiments reported later in the empirical part.

The second objective is to examine the limits and generality of the original framework. While Ref.~\cite{Niu2019} demonstrates strong performance under Gaussian amplitude noise, it leaves open several broader questions about the nature of the learned robustness and about the comparison with more standard quantum optimal-control techniques. This thesis therefore extends the original study in several directions. In particular, it asks whether a controller trained under one stochastic perturbation model can generalize to qualitatively different perturbations, including quasi-static drift, colored noise, parameter mismatch, and decoherence. It also asks how the reinforcement-learning strategy compares with more traditional quantum optimal-control methods when the Hamiltonian, constraints, target family, and evaluation metrics are kept fixed. Finally, it investigates how the different weights in the UFO objective influence the tradeoff between fidelity, leakage, boundary behavior, and runtime. These questions define the scope of the thesis beyond reproduction: the goal is not merely to restate the original results, but to evaluate how robust, transferable, and interpretable the framework remains when tested outside its original domain of validation.

It is also important to state what the thesis does \emph{not} attempt to do. The purpose is not to provide a complete survey of all reinforcement-learning methods for quantum control, nor to derive a universally optimal alternative to the reference framework. Likewise, the thesis does not aim to perform a full experimental hardware implementation on a superconducting device. Instead, its contribution is more focused: it develops a physically grounded simulation-based assessment of the framework of Ref.~\cite{Niu2019}, verifies its main claims in a reproducible setting, and then probes its behavior under controlled extensions that are scientifically motivated by open questions on robustness and control generalization. In this sense, the thesis is positioned at the intersection of reproduction, validation, and targeted methodological extension.

\section{Research questions}
\label{subsec:intro_research_questions}
% %\textbf{TODO:} Da schematizzare un pelo
The general objectives outlined above can be formulated more precisely in terms of a small number of concrete research questions. 
\begin{itemize}
    \item The first question concerns \emph{reproducibility}: can the main physical and numerical claims of the original UFO framework be recovered within an independently implemented simulation pipeline based on the same effective Hamiltonian, target gate family, and control metrics? This question is fundamental because the scientific value of any proposed control framework depends not only on its conceptual elegance, but also on the reproducibility of its reported gains in gate speed and robustness. In the context of the present thesis, this question is addressed through the reconstruction of the original control setting and through the experiments devoted to reproducing gate-time comparisons, average-fidelity curves, and fidelity-variance behavior under stochastic control noise.

    \item A second question concerns \emph{robustness generalization}. The original paper optimizes control policies in the presence of Gaussian amplitude fluctuations and evaluates them under the same family of perturbations \cite{Niu2019}. This naturally raises the question of whether the resulting robustness is specific to that training distribution or whether it generalizes to qualitatively different perturbation mechanisms. Put differently, if a policy is trained to be robust against one stochastic error model, does it also remain reliable under quasi-static drifts, temporally correlated noise, Hamiltonian parameter mismatch, or simple decoherence processes? This question is important both physically and methodologically. Physically, real devices are affected by several non-identical error sources simultaneously. Methodologically, the answer helps determine whether noisy-environment RL is merely learning distribution-specific compensation or whether it is capturing a broader form of control regularity. This question motivates the out-of-distribution robustness experiments introduced later in the thesis.

    \item A third question concerns \emph{methodological comparison}. Reinforcement learning is not the only available approach to analog quantum control, and the original paper itself positions its method against simpler direct optimization baselines \cite{Niu2019}. It is therefore important to ask how the RL-based framework compares with more standard quantum optimal-control strategies, such as GRAPE- or GOAT-like approaches, when the comparison is made under identical physical assumptions. More precisely, if one fixes the Hamiltonian model, amplitude constraints, target gate family, and evaluation metrics, does the RL-based method provide a measurable advantage in nominal performance, robust performance, or runtime, and under which conditions? This question is central to the scientific interpretation of the framework, because it determines whether the observed gains should be attributed to the physics of the cost function, to the RL optimization strategy, or to some interaction between the two.

    \item A fourth question concerns \emph{objective-function sensitivity and interpretability}. The UFO objective is built as a weighted combination of gate infidelity, leakage penalties, boundary penalties, and runtime. While this scalarization is one of the strengths of the framework, it also raises a natural question: how sensitive are the learned controls to the relative choice of these weights? In particular, how do the coefficients associated with fidelity, leakage, boundary conditions, and runtime modify the qualitative behavior of the optimized pulse? Do they reveal smooth tradeoff curves, or do they induce sharp transitions between qualitatively different control regimes? This question is scientifically relevant because it connects the empirical performance of the controller to the structure of the underlying physical objective. It also bears directly on the interpretability of the framework, since a robust optimization method should ideally not depend in an opaque way on finely tuned scalar weights.
\end{itemize}

Taken together, these research questions define the logical trajectory of the thesis. The first establishes whether the original framework can be reconstructed reliably. The second probes whether the learned robustness extends beyond the perturbation model used during training. The third evaluates the distinct contribution of the reinforcement-learning approach relative to more traditional control methods. The fourth examines how the multi-objective structure of the control cost shapes the resulting solutions. The empirical chapters of the thesis are organized around these questions, so that each numerical study addresses a clearly defined scientific issue rather than serving as an isolated computational exercise.

\section{Main contributions}
\label{subsec:intro_main_contributions}

The main contributions of this thesis can be divided into three closely connected parts. The first contribution is a structured reconstruction of the physical and methodological framework introduced in Ref.~\cite{Niu2019}. Rather than treating the reference paper only as a source of final numerical claims, this thesis develops its logic step by step: from the multilevel superconducting-qubit model and the subspace structure of the gmon architecture, to the time-dependent Schrieffer--Wolff treatment of leakage, the construction of the UFO objective, and the continuous-action reinforcement-learning formulation used for control optimization. In this sense, the thesis provides not only a numerical reproduction attempt, but also a systematic theoretical exposition of the framework, organized so as to make explicit the physical role of each component and the way in which they fit together within a realistic analog quantum-control problem.

The second contribution is an empirical validation of the original framework through a dedicated reproduction study. In particular, the thesis implements a consistent simulation pipeline for the target gate family considered in Ref.~\cite{Niu2019} and uses it to reproduce the main classes of results reported in the paper, namely the comparison of analog gate times with synthesized-gate baselines, the dependence of average fidelity on stochastic control-noise strength, and the corresponding fidelity-variance behavior. This part of the work is important because it tests whether the main claims of the original framework remain observable under an independent implementation, thereby providing a concrete assessment of its reproducibility. The associated numerical studies also serve as the reference point against which all subsequent extensions are interpreted.

The third contribution is the extension of the original framework along directions that are scientifically relevant but not fully explored in the reference paper. These extensions include, first, an out-of-distribution robustness analysis in which a controller trained under Gaussian amplitude noise is evaluated under qualitatively different perturbation models, including quasi-static drift, colored noise, parameter mismatch, and decoherence; second, a comparative analysis between the reinforcement-learning-based controller and more standard quantum optimal-control baselines under matched Hamiltonian, target-family, and evaluation conditions; and third, an ablation study on the coefficients of the UFO cost function in order to clarify the tradeoffs induced by the different objective terms. Taken together, these extensions aim to move the discussion beyond the question of whether the original method works in its native setting, toward the broader question of how robust, transferable, and interpretable the framework is when viewed as a general strategy for analog quantum control.

A further contribution of the thesis lies in the interpretive perspective adopted throughout the work. The goal is not only to report whether one optimizer outperforms another on a given metric, but also to connect those outcomes to the underlying physics of the control problem. For this reason, the empirical results are consistently interpreted through the conceptual themes developed in the background chapters, including the distinction between compiled and analog gate realization, the multilevel origin of leakage, the speed--leakage tradeoff, and the multi-objective structure of realistic control costs. In this way, the thesis aims to contribute not only a set of numerical results, but also a clearer understanding of how reinforcement-learning-based control interacts with the structure of superconducting quantum hardware.

\section{Thesis outline}
\label{subsec:intro_thesis_outline}

The remainder of this thesis is organized as follows. Chapter~2 develops the theoretical background required for the study, including quantum computation as a controlled dynamical problem, the role of gate fidelity and runtime as performance metrics, the distinction between compiled universal gate synthesis and direct analog control, the multilevel nature of superconducting qubits, the notions of computational subspace and leakage, the speed--leakage tradeoff, and the main open-loop, closed-loop, and reinforcement-learning approaches relevant to quantum control. This chapter establishes the general physical and methodological language used throughout the rest of the thesis.

Chapter~3 introduces the paper-specific framework and its physical implementation on the gmon architecture. It presents the effective Hamiltonian and control channels, the computational and leakage subspaces in the notation of the reference paper, the motivation for the time-dependent Schrieffer--Wolff treatment of leakage, the resulting leakage bound, and the construction of the Universal cost Function control Optimization objective. The same chapter then formulates the control problem as a reinforcement-learning task, explains the use of trust-region policy optimization in a stochastic training environment, introduces the target gate family \(N(\alpha,\alpha,\gamma)\), describes the baseline comparison strategy of the original paper, and summarizes the main findings that the present work seeks to reproduce and extend.

Chapter~4 contains the empirical assessment of the framework. It first presents the numerical protocol and implementation choices used in this thesis, including the evaluation metrics and the selected hyperparameter regime. It then reports the reproduction experiments on representative targets and on the original benchmark figures. After this, the chapter develops the extension studies proposed in this work: out-of-distribution robustness tests, comparisons against standard quantum optimal-control baselines under matched assumptions, and an ablation analysis of the coefficients entering the UFO objective. Where appropriate, the numerical results are interpreted in light of the theoretical considerations developed in the previous chapters.

Finally, Chapter~5 collects the main conclusions of the thesis. It summarizes the extent to which the original framework has been successfully reproduced, discusses what the extension experiments reveal about the generality and limitations of the method, and outlines possible future directions, including broader robustness studies, alternative control objectives, and possible hybrid open-loop / closed-loop refinements. In this way, the thesis progresses from conceptual background, to paper-specific framework, to empirical validation and extension, with the aim of providing both a faithful reconstruction of the original method and a critical assessment of its broader significance.

\section{AI usage policy}
\label{subsec:intro_ai_usage_policy}
% %\textbf{TODO:} to be written ( mention terrence tao work on collaboration and ai cost of IDEA generation is now zero)

As noted by the world-renowned mathematician Terence Tao, the emergence of large language models has already begun to reshape the modern research pipeline, not only in writing and technical implementation, but also in the broader organization of exploratory work. In particular, AI systems have significantly reduced the cost of idea generation, making it possible to produce, inspect, and compare a much larger number of tentative directions than was previously practical. In this sense, their main impact is not necessarily to replace depth of reasoning, but rather to expand breadth of exploration, while shifting the main epistemic bottleneck from generation to verification, selection, and validation \cite{Tao2024MachineAssistedProof,Tao2026AIInterview}. This perspective is especially relevant to the present thesis, where AI-based tools were used as auxiliary instruments in both the implementation and writing workflow.

More specifically, AI assistance was used on the coding side to generate limited code snippets, support debugging, clarify the usage of specific Python libraries, and accelerate minor implementation or refactoring tasks. On the writing and research side, it was used to assist the discovery of potentially relevant literature through semantic search and query refinement, to identify related work, and to support the drafting, rephrasing, and stylistic revision of selected passages. In all cases, however, AI remained a supportive and non-autonomous tool. The mathematical derivations, physical interpretations, modeling assumptions, numerical methodology, experimental design, source verification, and final presentation of results were independently checked and validated by the author. No scientific claim, citation, or conclusion was included in this thesis solely on the basis of AI-generated output.